\documentclass{article}
\usepackage[T1]{fontenc}
\usepackage{lmodern}
\usepackage{graphicx}
\usepackage{amsmath,amssymb}
\usepackage[a4paper,margin=2cm]{geometry}
\usepackage[absolute,overlay]{textpos}
\setlength{\TPHorizModule}{1mm}
\setlength{\TPVertModule}{1mm}
\usepackage[indonesian]{babel}
\usepackage{hyperref}
\setlength{\emergencystretch}{2em}
% unit: Halaman 34

\begin{document}

% === Halaman 34 (llm) ===

\section*{Pembentukan Kunci Putaran}

\begin{itemize}
    \item Setiap putaran di dalam algoritma \textit{Rijndael} menggunakan kunci putaran atau \textit{round key}. Kunci putaran dibangkitkan dari kunci eksternal dari pengguna yang dinamakan \textit{cipher key} dan disimbolkan dengan peubah \textit{key}. Pembangkitan semua kunci putaran dilakukan oleh fungsi \textit{KeyExpansion()}.
    \item Pembangkitan kunci putaran di dalam Algoritma Rijndael tergolong rumit dan agak sukar diterangkan. Tinjau sebuah larik \textit{key} yang panjangnya 16 \textit{byte} (16 elemen) dan $Nr = 10$ putaran. Sepuluh kunci putaran akan disimpan di dalam larik \textit{rk}. Elemen awal \textit{key} langsung menjadi $rk[0]$. Elemen-elemen kunci lainnya akan disimpan di dalam $rk[1], rk[2], \dots, rk[10]$.
\end{itemize}

\end{document}
